\documentclass[conference]{IEEEtran}
\usepackage{cite}
\usepackage{amsmath,amssymb,amsfonts}
\usepackage{algorithmic}
\usepackage{graphicx}
\usepackage{textcomp}
\usepackage{xcolor}
\usepackage{booktabs}
\usepackage{url}
\usepackage{hyperref}
\usepackage{caption}
\usepackage{subcaption}
\usepackage{multirow}

\hypersetup{colorlinks=true, linkcolor=blue, citecolor=blue, urlcolor=blue}

% ── Document ──────────────────────────────────────────────────────────────────
\begin{document}

\title{AI-Driven Skin Cancer Triage: Binary Classification of Dermoscopic Images from the HAM10000 Dataset}

\author{\IEEEauthorblockN{Abhijeet Kumar Shah}
\IEEEauthorblockA{\textit{Enroll. ID: 230036} \\
\textit{B.Tech -- CS \& AI}}
\and
\IEEEauthorblockN{Aditya Raj Sharma}
\IEEEauthorblockA{\textit{Enroll. ID: 230123} \\
\textit{B.Tech -- CS \& AI}}
}

\maketitle

\begin{abstract}
Early detection of skin cancer reduces mortality rates significantly. In this project, we implement a diagnostic triage tool utilizing the HAM10000 dataset...
\end{abstract}

\begin{IEEEkeywords}
Skin Cancer, Binary Classification, ResNet18, Computer Vision, HAM10000
\end{IEEEkeywords}

\section{Introduction}
Melanoma is the most dangerous form of skin cancer. Due to the scarcity of dermatologists everywhere...

\section{Literature Review}
Computer-Aided Diagnosis (CAD) has been a vital part of dermoscopy...

\section{Dataset Description}
We used the HAM10000 dataset consisting of 10,015 images. We formulated the problem as binary classification (Cancer vs Non-Cancer)...

\section{Methodology}
\subsection{Data Preprocessing}
Metadata normalization and handling missing values using median imputation...

\subsection{Baseline Metadata Models}
Logistic Regression and Random Forest leveraging age, sex, and localization...

\subsection{Computer Vision Extraction}
HOG and Color Histograms...

\subsection{Deep Learning Classification}
ResNet18 pre-trained on ImageNet...

\section{Experiments and Results}
We tracked Recall strictly due to the cost of False Negatives...

\section{Conclusion}
Transfer learning drastically outperforms traditional metadata baselines...

\end{document}
